\documentclass[10pt,a4paper]{amsart}
\usepackage[margin=19mm]{geometry}
\usepackage{amsmath,amssymb,mathtools}
\usepackage[scr=rsfs,cal=euler]{mathalfa}
\usepackage{xfrac}
\usepackage[unicode,hidelinks,bookmarks=false]{hyperref}
\newcommand{\ie}{i.e.\ }
\newcommand{\cf}{cf.\ }
\newcommand{\resp}{respectively\ }
\newcommand{\Subsection}{\subsection}
\providecommand{\nobreakdash}{\nobreak\hbox{-}}
\setlength{\emergencystretch}{2em}
\multlinegap0pt
\usepackage{bm}
\usepackage{bbm}
\usepackage{dsfont}
\usepackage{xspace}
\usepackage{cancel}
\usepackage{mathtools}
\usepackage{cleveref}
\usepackage{amsthm}
\makeatletter
\@ifundefined{theorem}{\newtheorem{theorem}{Theorem}}{}
\@ifundefined{lemma}{\newtheorem{lemma}[theorem]{Lemma}}{}
\@ifundefined{proposition}{\newtheorem{proposition}[theorem]{Proposition}}{}
\makeatother
\DeclareMathOperator{\diag}{diag}
\newcommand{\Q}{\mathcal{Q}}
\newcommand{\adjoint}{*}
\newcommand{\eqdef}{\coloneqq}
\newcommand{\qi}{\mathbf{i}}
\newcommand{\qj}{\mathbf{j}}
\newcommand{\qk}{\mathbf{k}}
\newcommand{\set}[1]{\left\{ #1 \right\}}
\newcommand{\tH}{\widetilde{H}}
\title[Quaternionic Perfect Sequences and Hadamard Matrices]{Quaternionic Perfect Sequences and Hadamard Matrices}
\author{Aidan Bennett, Curtis Bright, Paul Colinot, Ashwin Nayak}
\date{Source: arXiv:2601.22337v1 (CC BY 4.0)}
\begin{document}
\maketitle

\noindent\emph{Source and licence.} Quaternionic Perfect Sequences and Hadamard Matrices. Version arXiv:2601.22337v1 (CC BY 4.0), distributed under the Creative Commons Attribution 4.0 International licence, as recorded in the arXiv licence metadata for this exact version and shown on the abstract page https://arxiv.org/abs/2601.22337v1. Full rights notice: licenses/arxiv-2601.22337-CC-BY-4.0.txt.\par\smallskip

\noindent\textit{Editorial scope.} Counted unit: the Proposition whose source label is \texttt{prop-o7-nonequiv} in the article (source0.tex lines 1701--1704, source label \texttt{prop-o7-nonequiv}), delivered together with the article's own complete original proof of it (source0.tex lines 1705--1747, one \texttt{proof} environment of 43 lines). No mathematics is written, repaired or re-derived: statement and proof are the article's own text line for line. Required context retained uncounted: lem-eq-cyclic-qhm (lines 1586--1589); eq-o7qhm (lines 1652--1664); eq-o7qhm-k10 (lines 1673--1685); eq-o7qhm-badl (lines 1687--1699). Every definition, display and label of those lines is retained unchanged. Omitted from this package and not redistributed: 1--1585, 1590--1651, 1665--1672, 1686--1686, 1700--1700, 1748--1912. Nothing inside a retained excerpt is omitted or altered: every retained body is delivered exactly as the article's lines are written, so no omission receipt of any kind accompanies this package. Citations: the retained excerpts carry no citation command at all, so no bibliography entry of the article is transcribed and no reference is redistributed. Reference adaptation: none was needed and none was made; every reference inside the delivered unit resolves inside it, so no unresolved reference or citation is produced. Printed slips kept literal: no printed word, symbol, hypothesis or number of the counted statement or of its proof was altered. Layout and class: the article's own class is \texttt{article} with packages including babel, fontenc, booktabs, amsmath, amsthm, amssymb, mathtools, physics, algorithm, algorithmic, url, enumerate, hyperref, xspace; the wrapper is the lane's pinned \texttt{amsart} editorial wrapper at 10pt on A4, which restates the theorem-like environments the retained text needs and adds the article's own macro definitions that the retained text needs (\texttt{\textbackslash Q}, \texttt{\textbackslash adjoint}, \texttt{\textbackslash diag}, \texttt{\textbackslash eqdef}, \texttt{\textbackslash qi}, \texttt{\textbackslash qj}, \texttt{\textbackslash qk}, \texttt{\textbackslash set}, \texttt{\textbackslash tH}), with extra wrapper packages (bm, bbm, dsfont, xspace, cancel, mathtools, cleveref). The delivered numbering is the unit's own. Assets: no retained span contains a figure environment, a graphics-inclusion command, a PDF-page inclusion, a table or a TikZ picture; no image file is copied into this package, the package declares no asset dependency and no image file is redistributed with it. Licence: version arXiv:2601.22337v1 of this article is distributed under the Creative Commons Attribution 4.0 International licence, as recorded in the arXiv licence metadata for this exact version and shown on the abstract page https://arxiv.org/abs/2601.22337v1. A full rights notice is delivered at licenses/arxiv-2601.22337-CC-BY-4.0.txt. Characters: every retained excerpt is pure ASCII, so the delivered engine, XeLaTeX with its default Latin Modern text font, reports no missing character.
\begin{lemma}
\label{lem-eq-cyclic-qhm}
Let\/~$H$ be a cyclic QHM of order\/~$n$ with first row\/~$h$. Let\/~$\tH$ denote the QHM obtained by dephasing the first column of\/~$H$, i.e., $\tH \eqdef \diag(c)^\adjoint H$, where\/~$c$ is the first column of\/~$H$. Any normalized QHM equivalent to\/~$H$ may be expressed as\/~$\alpha h_k Q_1 \tH \diag(g(k))^\adjoint Q_2 \alpha^*$, for some unit quaternion\/~$\alpha$, $k \in [0,n)$, and permutation matrices\/~$Q_1$, $Q_2$, where\/~$g(k)$ is the vector given by the cyclic permutation of the entries of\/~$h$ by\/~$k$ places to the left. 
\end{lemma}

\begin{equation}
\label{eq-o7qhm}
R \eqdef
\begin{bmatrix}
 1 & -\qj & -q \qj & -\qj & 1 & -1 & -1 \\
-1 & 1 & -\qj & - q \qj & -\qj & 1 & -1 \\
-1 & -1 & 1 & -\qj & - q \qj & -\qj & 1 \\
 1 & -1 & -1 & 1 & -\qj & - q \qj & -\qj \\
-\qj & 1 & -1 & -1 & 1 & -\qj & - q \qj \\
- q \qj & -\qj & 1 & -1 & -1 & 1 & -\qj \\
-\qj & - q \qj & -\qj & 1 & -1 & -1 & 1 
\end{bmatrix} .
\end{equation}

\begin{equation}
\label{eq-o7qhm-k10}
S \eqdef
\begin{bmatrix}
1 & 1 & 1 & 1 & 1 & 1 & 1 \\
1 & -\qk q \qj & -\qk q \qk & \qk q & \qk q \qk & \qk q \qj & (\qk q)^2 \\
1 & -q^* \qi & \qi & -\qk & \qk & -\qi & \qi q \\
1 & -1 & \qj q^* \qj & \qk & -\qj & -\qk & q \\
1 & \qi & -1 & -\qi q^* \qj & -\qk & \qk & \qk q \\
1 & -\qj & \qj & -1 & -\qi q^* \qi & \qi & -q \\ 
1 & -\qi & -\qj & \qi & -1 & \qj q^* \qk & -\qi q 
\end{bmatrix} ,
\end{equation}

\begin{equation}
\label{eq-o7qhm-badl}
T \eqdef
\begin{bmatrix}
1 & 1 & 1 & 1 & 1 & 1 & 1 \\
1 & \qi q^* & \qk q \qj & -1 & -\qi & \qk & \qi \\
1 & -q^* & -\qi & -\qi q \qi & -1 & -\qk & \qk \\
1 & -\qj q^* & -\qj & \qj & -\qi q \qk & -1 & -\qi \\
1 & q^* & \qj & \qk & -\qj & \qk q \qk & -1 \\
1 & -\qi q^* & \qi & -\qj & \qj & -\qi & q \qi \\
1 & (\qj q^*)^2 & \qj q^* \qk & \qj q^* \qj & -\qj q^* & -\qj q^* \qj & -\qj q^* \qk
\end{bmatrix} .
\end{equation}

\begin{proposition}
\label{prop-o7-nonequiv}
The QHM\/~$R$ in Eq.~\eqref{eq-o7qhm} is neither equivalent to the QHM\/~$S$ defined in Eq.~\eqref{eq-o7qhm-k10} nor to the QHM\/~$T$ defined in Eq.~\eqref{eq-o7qhm-badl}.
\end{proposition}

\begin{proof}
By \Cref{lem-eq-cyclic-qhm}, any normalized matrix equivalent to~$R$ has the form
\[
\alpha h_k Q_1 \tH \diag(g(k))^\adjoint Q_2 \alpha^*
\]
for~$\alpha$, $h$, $k$, $Q_1$, $\tH$, $g(k)$, and~$Q_2$ as in the lemma. Let~$L \eqdef h_k \tH \diag(g(k))^\adjoint $.
It suffices to prove that there is some column of~$L$ that does not correspond to any column of~$S$ or~$T$ (after the application of the permutations~$Q_1, Q_2$ and automorphism~$x \mapsto \alpha x \alpha^*$ as in the above expression).

Consider the second column of~$L$.
Depending on the value of~$k$, this column equals one of the columns of the table
\begin{equation}
\label{eq-2nd-column}
\begin{array}{c|cccccccccccccc}
k & 0 && 1 && 2 && 3 && 4 && 5 && 6 \\
\hline
h_k h_{0}^* h_{1} h_{1+k}^*  & 1 && -\qj q^* && -q \qj && -1 && \qj && -\qj && \qj \\
h_k h_{6}^* h_{0} h_{1+k}^*  & -\qj && -q^* && -q && \qj && 1 && -1 && 1 \\
h_k h_{5}^* h_{6} h_{1+k}^*  & \qj && q^* && q && -\qj && -1 && 1 && -1 \\
h_k h_{4}^* h_{5} h_{1+k}^*  & -\qj && -q^* && -q && \qj && 1 && -1 && 1 \\
h_k h_{3}^* h_{4} h_{1+k}^*  & -1 && \qj q^* && q \qj && 1 && -\qj && \qj && -\qj \\
h_k h_{2}^* h_{3} h_{1+k}^*  & \qj q^* && (q^*)^2 && 1 && -q^* \qj && \qj q^* \qj && -\qj q^* \qj && \qj q^* \qj \\
h_k h_{1}^* h_{2} h_{1+k}^*  & \qj q  && 1 && q^2 && -q \qj && \qj q \qj && -\qj q \qj && \qj q \qj 
\end{array},
\end{equation}
since the~$i$th entry of the column for~$k \in [0,7)$ is~$h_k h_{-i}^* h_{1-i} h_{1+k}^*$.
The column is not all~$1$, and has one of the following types depending on the value of~$k$. 
\begin{enumerate}
\item \label{item-p1}
It has two elements in~$\set{\pm 1}$ and three elements in~$\set{ \pm \beta}$, where~$\beta \in \Q_8$ is some pure imaginary quaternion.

\item \label{item-p2}
It has three elements in~$\set{\pm 1}$ and two elements in~$\set{ \pm \beta}$, where~$\beta \in \Q_8$ is some pure imaginary quaternion.

\item \label{item-p3}
It has two distinct pairs of quaternions of the form~$(\beta, -\beta)$ for some~$\beta \in \Q_{24} \setminus \Q_8$, and one pair of quaternions~$(\gamma, \gamma^2)$ where~$\gamma \in \Q_{24} \setminus \Q_8$ is a cube root of~$-1$. 

\end{enumerate}

All elements of~$R$, $S$, $T$ are in~$\Q_{24}$, but the rest of the proof only uses the property that the elements of~$ \Q_{24} \setminus \Q_8$ have non-zero real and imaginary parts. Suppose~$R$ is equivalent to~$S$ or~$T$. Then the automorphism~$f \colon x \mapsto \alpha x \alpha^* $ (and the permutations~$Q_1$, $Q_2$) preserves the properties of the three types of columns described above. However, none of the columns of~$S$ or~$T$ are of types~\ref{item-p1} or~\ref{item-p2}. Only the last column of~$S$ and the second column of~$T$ are of type~\ref{item-p3}. So the second column of~$L$ may only be mapped to these columns, and with~$k \in \set{1,2}$ (cf. Eq.~\eqref{eq-2nd-column}). We argue next that this too is not possible.

We present the argument for the last column of~$S$ and~$k = 1$. The other cases follow along the same lines.
The only pairs of the form~$(\beta, -\beta)$ in the last column of~$S$ are~$(q,-q)$ and~$(\qi q, -\qi q)$, and in the second column of~$L$ with~$k = 1$ are~$(q^*, -q^*)$ and~$(\qj q^*, -\qj q^*)$. Since~$\Re(q) = \Re(-\qi q) = \Re(q^*) = \Re(\qj q^*) = 1/2$, the equivalence operations may only map~$(q^* \mapsto q, ~ \qj q^* \mapsto -\qi q)$ or~$(q^* \mapsto -\qi q, ~ \qj q^* \mapsto q)$. Further, we have~$-q^* \mapsto \qk q$ or~$-q^* \mapsto (\qk q)^2$. However, the two sets of conditions are not consistent as they map~$q^*$ to distinct quaternions, so the required equivalence operations do not exist.
\end{proof}

\end{document}
